\section{Lower bound}\label{sec:matching-lower}
Throughout the proof, $h=B^{-1/2}$ is the budget scale and
$\mathfrak h(P,Q)$ is the squared Hellinger distance defined in
Section~\ref{sec:notation}.

\subsection{An explicit posterior-martingale bound for the quantile tail}
Fix a gap $d$ and a capped one-coin policy. Under the equal sign prior,
let $U_k=\E[S\mid\text{seed},X_1,\ldots,X_k]$, with $X_i\in\{-1,1\}$
and $\E[X_i\mid S]=2dS$. Then $U_0=0$ and $U_k$ is a martingale.
For a current posterior mean $u$,
\[
 U_{k+1}=\frac{u+2dX_{k+1}}{1+2duX_{k+1}},\qquad
 \E[(U_{k+1}-u)^2\mid U_k=u]
 =\frac{4d^2(1-u^2)^2}{1-4d^2u^2}\le4d^2.
\]
At $d=1/2$, the posterior becomes absorbing at $\pm1$ after the first
observation; the undefined zero-over-zero expression there is interpreted
through the zero increment. For the policy's capped stopping time $N$,
\begin{equation}\label{eq:posterior-energy}
                         \E_{1/2}U_N^2\le4d^2\E_{1/2}N.
\end{equation}
The seed may be arbitrary and independent of the latent state.

Change the conditional positive-sign probability to
$1/2-\epsilon$ or $1/2+\epsilon$, $0\le\epsilon\le1/4$.
Relative to the midpoint transcript law the endpoint densities are
$1-2\epsilon U_N$ and $1+2\epsilon U_N$.
Since $1-\sqrt{1-x}\le x$ on $[0,1]$, (\ref{eq:posterior-energy}) gives
\begin{equation}\label{eq:explicit-tail-local}
 \mathfrak h(P_{1/2-\epsilon},P_{1/2+\epsilon})
 \le4\epsilon^2\E_{1/2}U_N^2
 \le16\epsilon^2d^2\E_{1/2}N.
\end{equation}

\begin{lemma}[Quantile-tail obstruction]\label{lem:tail}
For every allowed law, rarity and budget, $R_{\mu,\rho}(B)\ge T_h/32$.
\end{lemma}
\begin{proof}
Use an oracle experiment that lets the algorithm choose any tail atom
$i$ as an observable type, then obtain a genuine coin with that fixed
gap $d_i$ and unknown sign. Each flip is still charged. The algorithm
may interleave types and revisit coins.
Set
\[
 \epsilon_i=\frac1{16}\min\{1,h/d_i\},\qquad
 q_{i,\pm}=1/2\pm\epsilon_i.
\]
Equation (\ref{eq:explicit-tail-local}) bounds the local distance for
every type by $(h^2/16)$ times its midpoint work. The midpoint
decomposition, proved in Lemma~\ref{lem:tf-midpoint}, therefore bounds
global distance by $h^2B/16=1/16$. It permits observable types chosen
adaptively and arbitrary interleaving.

Every original algorithm can be simulated in this oracle experiment:
for each fresh coin privately choose tail type $i$ with probability
$\tau_i$; on the remaining mass simulate an all-negative coin with
the residual magnitude law $\kappa/(1-\rho)$ for free. Keep that
private choice on revisits. The resulting original endpoints are
\[
 q_\pm(d)=\frac{\mathrm d\nu}{\mathrm d\mu}(d)\left(\frac12\pm\frac1{16}\min\{1,h/d\}\right).
\]
The density uses total masses at a support point. At a distinct atom it
equals $\tau_i/w_i$; this remains well-defined even if an input lists
identical gaps more than once.
Both endpoints and their midpoint are admissible, since their total
positive mass is at most $(9/16)\rho$. The oracle's genuine work is
at most the original flip count, so its midpoint budget is at most $B$.
No budget assumption over arbitrary oracle parameters is needed.
The original target separation is
$2\sum_i\tau_i\epsilon_i=T_h/8$.
Data processing gives total variation at most
$\sqrt{2/16}<1/2$. Two-point midpoint testing gives RMS at least one
quarter of the separation, namely $T_h/32$.
For unbounded original stopping, cap the global experiment first.
The same endpoint pair works at every cap; increasing prefixes give
the terminal TV bound, as proved after Lemma~\ref{lem:tf-midpoint}.
\end{proof}

\begin{lemma}[Rare Bernoulli-label floor]\label{lem:floor}
$R_{\mu,\rho}(B)\ge\min\{\rho,h\sqrt\rho\}/32$.
\end{lemma}
\begin{proof}
Restrict to constant priors $q(d)=p$ and strengthen the experiment by
revealing a fresh coin's sign for unit cost. Its magnitude and subsequent
flips can then be simulated using independent randomness, so the
number of purchased labels is at most the original number of flips.
Take $p_- =\rho/2$, $p_+=p_-+\delta$ with
$\delta=\min\{\rho,h\sqrt\rho\}/8$.
One Bernoulli-label divergence is at most
$\delta^2/[p_+(1-p_+)]\le3\delta^2/\rho$.
The stopped chain rule bounds total KL by $3\delta^2B/\rho\le3/64$.
Pinsker gives TV below $1/2$, and midpoint testing gives RMS
at least $\delta/4$. Finite prefixes justify unbounded stopping.
\end{proof}

\subsection{A deterministic entropy filter}
Retain the canonical core notation $a,A,F,\lambda$ from
Section~\ref{sec:entropy-preliminaries}.
If $h\ge t$, then $T_h=\rho$, so the tail lemma already proves the
desired lower bound. Assume $h<t$. If $\lambda\ge t/4$, then
\[
 h\sqrt{\rho\Gamma_\kappa}\le2h\sqrt{\rho F}
 =2h\rho/\lambda\le8h\rho/t\le8T_h,
\]
and $\Q\le9T_h$. Thus it remains to treat $\lambda<t/4$.

Set $L_j=1+\log(A/a_j)$ on positive-cost bands. Declare a band
\emph{good} exactly when
\begin{equation}\label{eq:good-band}
                       b_j\ge\tfrac\lambda2 L_j.
\end{equation}
This is a direct test of the known weights and endpoints.
Let $F_g$ and $F_b$ be the entropies of the good and bad cost vectors,
each normalized by its own total cost.

\begin{lemma}[Filtering retains core entropy]\label{lem:filter}
If $\lambda<t/4$, then $F_b\le F/8$ and $F_g\ge F/4$.
\end{lemma}
\begin{proof}
Every core band has endpoint at least $t$. Number the bands upward
from the one containing $t$, with endpoints $\beta_0 2^k$, $k\ge0$.
On a bad band, $L_j>x_j:=2b_j/\lambda$, and $x_j\ge8$.
Since $a_j/A=e^{1-L_j}$ and $xe^{1-x}$ decreases for $x\ge1$,
\[
 \sum_{j\text{ bad}}a_jL_j
 \le A\sum_{k\ge0}x_0 2^k e^{1-x_0 2^k},\qquad x_0=2\beta_0/\lambda\ge8.
\]
The ratio of successive summands is at most $2e^{-8}<1/2$.
Their sum is at most $2x_0e^{1-x_0}\le16e^{-7}<1/8$.
The bad vector's own entropy is at most this full-normalization sum,
so $F_b\le A/8\le F/8$. By (\ref{eq:binary-mixture}),
$F\le2(F_g+F_b)$; consequently $F_g\ge3F/8\ge F/4$.
\end{proof}

\subsection{Packing the filtered entropy into actual experiments}
We first group coordinates with comparable costs.
\begin{lemma}[Entropy tiers]\label{lem:entropy-tiers}
For a finite positive vector $a$, let $M=\max_j a_j$ and group its
coordinates by $a_j\in(M2^{-k-1},M2^{-k}]$, $k\ge0$.
If $F_k$ is each group's entropy cost, then
\begin{equation}\label{eq:entropy-tiers}
 \mathcal G(a)\le7\sum_k F_k.
\end{equation}
\end{lemma}

\begin{proof}
Let $A_k$ be the cost mass in group $k$, $n_k$ its cardinality, and
$p_k=A_k/A$. Entropy decomposition gives
$\mathcal G(a)=\sum_kF_k+A\Ent(p)$.
Empty groups are omitted. For the comparison on $k\ge0$, the reference
probability is $q_k=2^{-k/2}/(2+\sqrt2)$, so $\KL(p\Vert q)\ge0$
applies even when the nonempty tier indices are sparse.
Comparison with the probability vector proportional to $2^{-k/2}$ gives
\[
 A\Ent(p)\le A\log(2+\sqrt2)
                 +\frac{\log2}{2}\sum_k kA_k.
\]
On groups with $n_k\ge2^{k/2}$, the second term is at most
$\sum_kA_k\log n_k\le\sum_kF_k$; comparable weights imply
$F_k\ge A_k[1+\log n_k-\log2]\ge A_k\log n_k$.
On the other groups, $A_k\le M2^{-k/2}$, so their contribution is at most
\[
 \frac{\log2}{2}M\sum_{k\ge0}k2^{-k/2}
 =\frac{(\log2)\sqrt2}{2(\sqrt2-1)^2}M<3M\le3A.
 \]
Finally $\log(2+\sqrt2)<2$ and $A\le\sum_kF_k$.
Adding the initial $\sum_kF_k$ proves \eqref{eq:entropy-tiers}.
\end{proof}


Split the good bands into high bands $j<6$ and low bands $j\ge6$,
with entropies $F_H,F_L$. Since $b_j\ge1/64$ on the first six bands,
$F_H\le4096(1+\log6)=:C_H$.
If $F_H\ge F_L$, then
\[
 F\le4F_g\le8(F_H+F_L)\le16C_H,
\]
and the Bernoulli-label floor implies
$R\gtrsim\min\{\rho,h\sqrt{\rho F}\}$.

Otherwise $F_L>F_H$, so $F\le16F_L$.
Split the low bands by index parity, then within each parity into the
comparable-cost tiers of Lemma~\ref{lem:entropy-tiers}.
The sum of the resulting group entropies is at least $F_L/14$.
Indeed $F_L\le2(F_{\rm even}+F_{\rm odd})$, and the tier lemma is
applied separately within each parity. Both parities and all nonempty
tiers are retained; choosing just one tier would not justify this bound.
Each group has endpoint ratios at least four and costs in $(c,2c]$
for some $c>0$.

For a group of $n$ bands, (\ref{eq:good-band}) and $A\ge nc$ give
\[
 b_j\ge\frac\lambda2[1+\log(n/2)]
       \ge\frac\lambda8(1+\log n).
\]
Remove its $s=\lceil\log_4n\rceil$ smallest-gap bands, with $s=0$
for $n=1$. For $n\ge2$, $s\le n/2$, so at least $n/2$ remain.
Let $G_i$ be the remaining group's entropy. Coordinate monotonicity yields
\begin{equation}\label{eq:trimmed-mass}
 G_i\ge F_{\rm group}/16,\qquad
 b_{\min,i}\ge\lambda n(1+\log n)/8.
\end{equation}
For the first bound, divide the lower bound
$cn'(1+\log n')$ by the upper bound $2cn(1+\log n)$;
the ratio is at least $(1-\log2)/4>1/16$.
The second bound follows because deleting $s$ smallest bands raises
the minimum endpoint by at least $4^s\ge n$.
The case $n=1$ is immediate. The integer inequality for $n\ge2$
can be checked directly at $n=2,3,4$; when
$4^{s-1}<n\le4^s$ and $s\ge2$, we have
$2s\le4^{s-1}+1\le n$. Trimming is by increasing gap (equivalently,
increasing endpoint), not by cost. The number $n$ in
\eqref{eq:trimmed-mass} is the pre-trim cardinality.

Let $\sigma_i$ be the restriction of the actual core measure to a
remaining group, $m_i$ its mass, $u_i$ its largest endpoint, and
$w_{*,i}$ its smallest band mass. Then
\[
 m_i/u_i^2\ge c,\qquad w_{*,i}\ge c b_{\min,i}^2,
 \qquad G_i\le2cn(1+\log n).
\]
Consequently
\begin{equation}\label{eq:direct-saturation}
                 \frac{m_iw_{*,i}}{u_i^2}\ge\lambda^2G_i^2/256.
\end{equation}
Assign original rarity $\theta_i=\lambda^2G_i/1024$ and conditional
rarity $r_i=\theta_i/m_i$.
Equation (\ref{eq:direct-saturation}) gives
$\theta_iG_i\le m_iw_{*,i}/(4u_i^2)$.
Because $G_i\ge m_i/u_i^2$, this actually implies
$\theta_i\le w_{*,i}/4\le m_i/4$, hence $r_i\le1/4$.
Here $w_{*,i}$ is a band mass, not a minimum individual atom mass.
Under the conditional law $\sigma_i/m_i$, the entropy and minimum
band mass are $G_i/m_i$ and $w_{*,i}/m_i$. Its entropy-branch condition is
\[
 r_i u_i^2(G_i/m_i)\le w_{*,i}/m_i
 \quad\Longleftrightarrow\quad
 \theta_iG_i\le m_iw_{*,i}/u_i^2.
\]
Thus the preceding factor-four margin verifies the actual conditional
condition. There is no unnormalized-to-conditional mass factor missing.
Disjointness and superadditivity give
\[
 \sum_i\theta_i\le\lambda^2F/1024=\rho/1024,
 \qquad \sum_iG_i\ge F_L/224\ge F/3584.
\]

All remaining bands have largest endpoint at most $1/128$ and
endpoint ratios at least four. Apply the annular one-unit obstruction
of Proposition~\ref{prop:scaled-annular} to each normalized component
$\sigma_i/m_i$, followed by the additive localization theorem
(Theorem~\ref{thm:additive-localization}). Their complete proofs are
included in Appendices~\ref{sec:lower-foundations} and
\ref{sec:transfer-foundations}.
The actual minimum-band-mass condition is verified by
(\ref{eq:direct-saturation}), so the entropy branch is active.
We obtain
\[
 R(B)\gtrsim\sum_i\min\{\theta_i,h\sqrt{\theta_iG_i}\}
 =\left(\sum_iG_i\right)
       \min\{\lambda^2/1024,h\lambda/32\}
 \gtrsim\min\{\rho,h\sqrt{\rho F}\}.
\]
The proportional allocation $\theta_i/G_i=\lambda^2/1024$ is what
makes this estimate valid jointly at every budget.

\subsection{Completion of the lower bound}
The cases $h\ge t$ and $\lambda\ge t/4$ were handled by the tail lemma.
In the remaining case, the floor or the packing argument gives
$R\gtrsim\min\{\rho,h\sqrt{\rho F}\}$, while always $R\ge T_h/32$.
Finally, by (\ref{eq:canonical-core}),
\[
 \Q(B)\le T_h+2\min\{\rho,h\sqrt{\rho F}\}.
\]
The two lower bounds control this sum up to a universal constant.
Thus $R(B)\ge c\Q(B)$, proving the lower half of
Theorem~\ref{thm:direct-main}. Section~\ref{sec:assembly-contract}
tracks the constants to obtain the stated value $10^{-20}$.
